\section{Розділ~\ref{sec:synthesis}. Уся машина}

Розділ-зведення нових дослідів не наводить: кожен його рядок спирається на розділ, де твердження розібрано, і на ті самі джерела; твердо встановлено шину \nt{GLUT} і керування потоком через \nt{GABA}, а ролі широкомовних каналів та їхня спільна робота --- здебільшого гіпотеза.

{\footnotesize
\setlength{\tabcolsep}{3pt}\renewcommand{\arraystretch}{1.15}
\begin{longtable}{>{\raggedright}p{0.5cm}>{\raggedright}p{4.0cm}>{\raggedright}p{5.6cm}>{\raggedright}p{2.3cm}>{\raggedright\arraybackslash}p{1.7cm}}
\toprule
№ & Твердження & Дослід і що виміряно & Джерело & Статус \\
\midrule\endfirsthead
\toprule
№ & Твердження & Дослід і що виміряно & Джерело & Статус \\
\midrule\endhead
1 & На $8.6\cdot10^{10}$ нейронів припадає лише чотири види керувальних сигналів~\eqref{eq:machine} & Підрахунок ядер клітин у гомогенаті мозку (ізотропний фракціонатор): у дорослого чоловіка $86.1\pm8.1$ млрд нейронів & \cite{Azevedo2009} & виміряно \\
2 & Твердження~\ref{prop:architecture}, перші два шари: дані йдуть адресною шиною \nt{GLUT}, знак зв'язку задає відправник, потік спрямовують і нормують \nt{GABA}-нейрони (розділи~\ref{sec:neurontypes}, \ref{sec:gaba}) & Один головний медіатор на нейрон (огляд вимірювань); пряме гальмування звужує вікно, в якому пірамідний нейрон додає входи; ділення на загальну активність виміряно в багатьох ділянках & \cite{Strata1999,Pouille2001,Carandini2012} & виміряно \\
3 & Елементи ненадійні: синапс губить більшість імпульсів (табл.~\ref{tab:computer}, розділ~\ref{sec:code}) & Зрізи гіпокампа, мінімальна стимуляція: понад половина імпульсів не викликає відповіді в CA1; збої порога й проведення аксоном виключено, причина --- імовірнісне вивільнення медіатора & \cite{AllenStevens1994} & виміряно \\
4 & Мозок працює на $\sim20$ ватах, у середньому близько імпульсу за секунду на нейрон (розділ~\ref{sec:code}); такої машини інженери поки не збудували & Оцінка~\eqref{eq:energy} з виміряних витрат: близько $10^9$ молекул АТФ на імпульс разом із роботою синапсів (кора гризунів); детальна оцінка для кори дає ще меншу частоту. Стан техніки --- за оглядом & \cite{Attwell2001,Lennie2003,MehonicKenyon2022} & теорія \\
5 & Навчання більшості синапсів --- добуток місцевого сліду й одного спільного числа $\delta$ (друге рівняння~\eqref{eq:machine}, розділ~\ref{sec:dofa}) & \nt{DOPA}-нейрони мавпи відповідають на помилку передбачення винагороди; у зрізах смугастого тіла дофамін збільшує шипик, лише якщо надходить через $0.3$--$2$\,с після глутаматного входу --- слід виміряно & \cite{Schultz1997,Yagishita2014} & виміряно \\
6 & Окремий провід помилки до кожного виходу є лише в колі навчання рухів (розділ~\ref{sec:motorlearning}) & Мозочок: збіг сигналу ліаноподібного волокна з входом послаблює цей вхід; у мавпи складні імпульси клітин Пуркіньє несуть напрямок помилки руху й викликають поправку & \cite{Ito2001,Herzfeld2018} & виміряно \\
7 & Кожен широкомовний канал --- джгут із кількох ліній (твердження~\ref{prop:bundle}) & Для \nt{DOPA}: в одній задачі різні нейрони несуть винагороду, рух і ознаки стимулу; швидка помилка й повільний рівень дофаміну різні. Для \nt{SERO}, \nt{NORA}, \nt{ACET} такий розклад вивчено гірше & \cite{Engelhard2019,Hamid2016,Mohebi2019} & виміряно \\
8 & \nt{SERO}: регулятор рівня і, за гіпотезою, горизонт $\tau_\gamma$ (розділ~\ref{sec:serocomputer}) & Самогальмування: агоністи власних рецепторів 5-HT$_{1A}$ гальмують серотонінові нейрони шва --- виміряно. Горизонт запропоновано в теорії; найсильніший дослід --- оптогенетична активація цих нейронів у миші подовжує очікування відкладеної винагороди & \cite{Sprouse1987,Doya2002,Miyazaki2014} & гіпотеза \\
9 & \nt{NORA}: підсилення $g$ вибирає між пошуком і рішенням (розділ~\ref{sec:nora}) & Зростання відношення сигнал/шум: у мавпи в стані неспання норадреналін у слуховій корі пригнічує фонову активність сильніше, ніж відповідь на голоси родичів --- виміряно; мультиплікативне підсилення --- модель; роль у виборі між пошуком і рішенням --- теорія & \cite{Foote1975NE,ServanSchreiber1990,AstonJones2005} & гіпотеза \\
10 & \nt{ACET}: перемикає запис і читання (розділ~\ref{sec:acet}) & Три дії (послаблення внутрішніх зв'язків, підсилення входу, полегшення пластичності) виміряно; перемикання режимів опосередковано підтримане дослідом зі скополаміном у людини & \cite{HasselmoBower1992,Gil1997,Atri2004} & гіпотеза \\
11 & Формула~\eqref{eq:machine} описує машину цілком & Це зведення моделей розділів~\ref{sec:glutcomputer}--\ref{sec:acet}, кожна частина спирається на свій розділ; як ціле дослідом не перевірялася & виведення в розділі & гіпотеза \\
12 & Ручки працюють злагоджено без спільного керування: помилка, несподіванка, запис, повернення (рис.~\ref{fig:timeline}) & Досліду немає: схема якісна. Окремі ланки --- теорія розподілу праці \nt{NORA} і \nt{ACET} та зв'язок зіниці зі швидкістю навчання в людини & \cite{YuDayan2005,Nassar2012} & гіпотеза \\
13 & Навчання за одним спільним сигналом тим повільніше, що більше нейронів впливає на результат (твердження~\ref{prop:broadcastcost}) & Розрахунок кривих навчання лінійної мережі: навчання за збуренням виходів повільніше за градієнтний спуск у стільки разів, скільки виходів & \cite{Werfel2005} & теорія \\
14 & Як кора налаштовує багатошарові кола, невідомо; кандидати --- пачки імпульсів, обчислення в гілках нейрона, самонавчання передбаченням & Пачки як носій помилки --- модель; повторити відповідь детальної моделі нейрона кори змогла лише мережа з $5$--$8$ шарів --- модель; кальцієві імпульси в гілках нейронів кори людини, що дають виключне АБО, --- виміряно в зрізах; самонавчання --- огляд свідчень & \cite{Payeur2021,Beniaguev2021,Gidon2020,Keller2018} & гіпотеза \\
\bottomrule
\end{longtable}}
